\documentclass[11pt,letterpaper,serif]{moderncv}

\usepackage{amsmath}
\usepackage{amssymb}
\usepackage{hyperref}
\hypersetup{
	colorlinks=true,
	linkcolor=blue,
	filecolor=green,      
	urlcolor=green
}
\moderncvstyle{banking} % 'casual', 'classic', 'oldstyle' 'banking', 'fancy'
\moderncvcolor{red} % 'blue' (default), 'orange', 'green', 'red', 'purple', 'grey' and 'black'

\usepackage[scale=0.75]{geometry} % Reduce document margins
\addtolength{\voffset}{-3.5em}
\addtolength{\textheight}{7em}
\setlength{\hintscolumnwidth}{2cm} % width of dates column

%\quote{``It's hard to do a really good job
%on anything you don't think about in the shower.''
%--- Paul Graham}

\newcommand{\entry}[6]%
{\cventry[0.85em]{#4}{#3}{#2}{#1}{#5}{#6}} % banking
\newcommand{\cvurl}[1]{\bgroup\color{blue}\normalfont\url{#1}\egroup}
\newcommand{\arxiv}[1]{\bgroup\color{blue}\normalfont\rmfamily%
\href{http://arxiv.org/abs/#1}{arXiv:#1}\egroup}

%\usepackage[style=chem-rsc,maxnames=7]{biblatex}
%\addbibresource{biblio.bib}

\lfoot{\emph{Last updated \today}}

\newcommand{\subtext}[1]{\bgroup\footnotesize\rmfamily\color{black!70}#1\egroup}
\newcommand{\previous}[1]{\bgroup\color{black}\bfseries#1\egroup}

\newcommand{\job}[4]{\smallskip\cvitemwithcomment{#2}{#1}{#3}%
\unskip\subtext{#4}\medskip}

\newcommand{\pub}[2]{\cvlistitem{#1. \\ \subtext{#2}}\smallskip}
\newcommand{\paper}[3]{\pub{#1}{#2. \arxiv{#3}.}}

\name{Jalen}{Chrysos}
\email{jzchrysos11@gmail.com}
\social[github]{JZChrysos}
\lfoot{\emph{Last updated \today}}

\begin{document}

\makecvtitle % Print the CV title
\vspace*{-1.5em}

\section{Education}
\entry{Fall 2022-Spring 2026}
	{University of Chicago}%
	{BA/MS in mathematics.}
	{}{}{}
%\textbf{Relevant coursework}:
%\begin{itemize}
%	\item Honors Analysis in $\mathbb{R}^n$ 207-209.
%	\item Honors Basic Algebra 257-259.
%	\item 
%	\item Graduate Algebra I-III
%	\item Graduate Analysis I-III
%	\item Graduate Topology I-III
%\end{itemize}
\cvitemwithcomment{Reed College}%
	{Concurrent enrollment during high school.}{2020-2021}
\cvitemwithcomment{Portland State University}%
	{Dual enrollment during high school.}{2019-2021}

\subsection{Programming experience}
Python, C++, \LaTeX, Asymptote.

\section{Selected work experience}
\subsection{Research internships}

\vspace{3mm}

\job{Research Experience for Undergraduates}
{University of Chicago}{Summer 2025}
{Wrote an \href{https://github.com/JZChrysos/Math/blob/main/Reverse_Math_REU/Reverse_Math_REU.pdf}{expository paper} on Reverse Mathematics, with a focus on the relationships between Ramsey's Theorem, K\"onig's Lemma, and principles in real analysis. Extensive use and discussion of the method of forcing.}

\job{Algorithm Design Internship}
{Kepler Computing}{Summer 2023}
{Wrote algorithms in python for optimal synthesis of logic functions from majority gates.}

\job{Software Internship}
{Saltire Software}{Summer 2019}
{Worked on a team to adapt Hermann Helmholtz's \textit{On the Sensations of Tone} into an \href{https://sensationsoftone.com/}{interactive e-book}. Wrote audio-physics applets in Java and typeset diagrams in \LaTeX.}

\job{Data Analyst and Electrical Testing Engineer}
{Kepler Computing}{June 2021 - September 2022}
{Wrote testing automation code in python and performed data analysis using matplotlib for the design of ``Beyond CMOS'' computing hardware. Employed full-time.}

\subsection{Teaching and Mentorship}

\vspace{3mm}

\job{Teaching Assistant}
{University of Chicago}{2023}
{TA for 160s IBL (Inquiry Based Learning) sequence, an introduction to proofs in real analysis. Graded problem sets, held weekly office hours.}

\job{Math Contest Tutor}
{Private Tutoring}{2023-present}
{Taught kids olympiad-level material in a creative and fun way.}

\job{Team Leader}{Oregon ARML}{2020-2021}{Team leader representing Oregon at the ARML (American Regional Math League) competition. Organized practices and logistics for competing, mentored younger kids on the team.}

\subsection{Other}

\vspace{3mm}

\job{Projectionist}{Doc Films}{2023-present}
{Worked weekly shifts as an apprentice projectionist. Projected and prepped 35mm and 16mm film, as well as digital.}

\section{Selected Fellowships and Achievements}

\begin{itemize}
	\item \textbf{Metcalf Grant recipient}: For work at the UChicago Math REU (Summer 2025)
	\item \textbf{Harvard MIT Math Tournament}: Sixth place team (February 2020)
	\item \textbf{USA Junior Math Olympiad participant}: Rank $\sim$70 nationally (Spring 2018)
	\item \textbf{North American Linguistics Olympiad}: Semifinalist and Best Solution Award (Spring 2021)
\end{itemize}



\end{document}
